% 大容量假设空间、插值解集合与优化算法的选解偏差。
\begin{tikzpicture}[x=1mm,y=1mm,font=\footnotesize]
  \fill[FigureInput!7] (0,0) ellipse (49 and 32);
  \draw[FigureInput,thick] (0,0) ellipse (49 and 32);
  \node[font=\kaishu\scriptsize,text=FigureInput] at (0,27) {架构可表示的假设空间};

  \draw[FigureLoss,very thick]
    plot[smooth cycle,tension=0.8] coordinates {(-31,-7) (-15,10) (5,13) (28,4) (17,-14) (-8,-18)};
  \node[font=\kaishu\scriptsize,text=FigureLoss] at (-3,-23) {经验风险为零的插值解};

  \draw[BookMuted,dashed,->] (-45,25) .. controls (-31,19) and (-27,8) .. (-18,4);
  \draw[BookMuted,dashed,->] (-45,18) .. controls (-25,11) and (-12,5) .. (0,2);
  \draw[BookMuted,dashed,->] (-45,11) .. controls (-22,4) and (6,0) .. (17,-3);
  \fill[FigureModel] (0,2) circle (2.2);
  \node[font=\kaishu\scriptsize,text=FigureModel,align=left,anchor=west] at (5,16)
    {优化器、初始化与正则化\\共同选择的一个解};
  \draw[book arrow,FigureModel] (7,14) -- (1.5,4);

  \node[font=\kaishu\scriptsize,text=BookMuted,align=center,text width=95mm] at (0,-38)
    {容量描述整个假设空间；泛化还取决于算法选择了哪个假设。};
\end{tikzpicture}
